%%
%% ACV 2026 submission
%% Based on the PLanQC 2025 extended-abstract template.
%% Compile: lualatex main; bibtex main; lualatex main; lualatex main.
%%
\documentclass[sigplan, nonacm]{acmart}

\AtBeginDocument{%
  \providecommand\BibTeX{{%
    Bib\TeX}}}

% \acmConference[ACV 2026]{First Workshop on Abstract and Concrete Techniques in Verification}{July 24--25, 2026}{Lisbon, Portugal}
% \acmBooktitle{ACV 2026: First Workshop on Abstract and Concrete Techniques in Verification, July 24--25, 2026, Lisbon, Portugal}


% Remove the permission/copyright text
% \settopmatter{printacmref=false} % Removes the reference format block
% \renewcommand\footnotetextcopyrightpermission[1]{}
% \pagestyle{plain} % Removes running headers

% MY OWN THINGS

% \usepackage{amsmath,amssymb,amsthm}
% \usepackage{amsthm}
% \usepackage[margin=2cm]{geometry}
\usepackage{dsfont}
\usepackage{tikz}
\usetikzlibrary{cd}
\usetikzlibrary{babel}
% \usepackage{tikz-cd}
\usepackage{titlesec}
\usepackage{hyperref}
\usepackage{cleveref}
\usepackage{parskip}
\usepackage{ebproof}
% \usepackage{quiver}
\usepackage{stackengine,scalerel}
% \usepackage{xparse}
% \usepackage{physics}
\usepackage{physics2}
\usephysicsmodule{ab,ab.braket}
\usepackage{tikzit}
\input{smc.tikzstyles}

\usepackage{mathtools}
\usepackage{stmaryrd}

% \usepackage{unicode-math}
% \setromanfont{TeX Gyre Termes}
% \setsansfont{TeX Gyre Heros}
% \setmonofont{TeX Gyre Cursor}[Ligatures=NoCommon]
% \setmathfont{TeX Gyre Termes Math}

% \usepackage{luatexja}

\usepackage{listings}% http://ctan.org/pkg/listings
\lstset{
  basicstyle=\ttfamily,
  mathescape,
  language=Python,
  morekeywords={qbit, bit, out, in, new, measure, apply, skip}
}

% \usepackage{quantikz}


%% Theorems
\theoremstyle{plain} % default

\newtheorem{theorem}{Theorem}[section]
\newtheorem{proposition}[theorem]{Proposition}
\newtheorem{lemma}[theorem]{Lemma}
\newtheorem{claim}[theorem]{Claim}
\newtheorem{corollary}[theorem]{Corollary}
\newtheorem{remark}[theorem]{Remark}
\newtheorem{definition}[theorem]{Definition}
\newtheorem{convention}[theorem]{Convention}
\newtheorem{conjecture}[theorem]{Conjecture}
\newtheorem{example}[theorem]{Example}
\newtheorem{nonexample}[theorem]{Non-Example}




%% Probability
\makeatletter
\renewcommand{\P}{%
	\@ifnextchar\bgroup%
	{\@Pwithargs}
	{\@Pnoargs}
}
\newcommand{\@Pwithargs}[1]{%
	\@ifnextchar\bgroup%
	{\@Ptwoargs{#1}}
	{\@Ponearg{#1}}
}
\newcommand{\@Pnoargs}{\mathbb{P}}
\newcommand{\@Ponearg}[1]{\mathbb{P}\left[ #1 \right]}
\newcommand{\@Ptwoargs}[2]{\mathbb{P}_{#1}\left[ #2 \right]}
\newcommand{\E}{%
	\@ifnextchar\bgroup%
	{\@Ewithargs}
	{\@Enoargs}
}
\newcommand{\@Ewithargs}[1]{%
	\@ifnextchar\bgroup%
	{\@Etwoargs{#1}}
	{\@Eonearg{#1}}
}
\newcommand{\@Enoargs}{\mathbb{E}}
\newcommand{\@Eonearg}[1]{\mathbb{E}\left[ #1 \right]}
\newcommand{\@Etwoargs}[2]{\underset{#1}{\mathbb{E}}\left[ #2 \right]}
\makeatother

%% MACROS
\newcommand{\naturals}{\mathbb{N}}
\newcommand{\integers}{\mathbb{Z}}
\newcommand{\indicator}{\mathds{1}}
\newcommand{\va}{\mathbf{a}}
\newcommand{\vb}{\mathbf{b}}
\newcommand{\vx}{\mathbf{x}}
\newcommand{\vy}{\mathbf{y}}
\newcommand{\vw}{\mathbf{w}}
\newcommand{\vz}{\mathbf{z}}
\newcommand{\vsigma}{\mathbf{\sigma}}
\newcommand{\reals}{\mathbb{R}}
\newcommand{\complex}{\mathbb{C}}
\newcommand{\field}{\mathbb{F}}
\newcommand{\sign}{\mathrm{sign}}
\newcommand{\lnorm}[2]{\Vert #1 \Vert_{#2}}
\newcommand{\ltwonorm}[1]{\lnorm{#1}{2}}
\newcommand{\grad}{\nabla}
\newcommand{\LinW}{\mathsf{Lin}_W}
\newcommand{\calF}{\mathcal{F}}
\newcommand{\calW}{\mathcal{W}}
\DeclareMathOperator*{\argmin}{arg\,min}
\newcommand{\PARTITION}{\mathsf{PARTITION}}
\newcommand{\NULLPLAYER}{\mathsf{NULLPLAYER}}
\newcommand{\KNAPSACK}{\mathsf{KNAPSACK}}
\newcommand{\Core}{\mathbf{Core}}
\newcommand{\ve}{\mathbf{e}}
\newcommand{\vzero}{\mathbf{0}}
\newcommand{\true}{\mathsf{true}}
\newcommand{\false}{\mathsf{false}}
\newcommand{\err}{\mathrm{err}}
\newcommand{\brackets}[1]{\{#1\}}
\newcommand{\polyinput}{\text{poly}(\text{input})}
\newcommand{\Ecal}{\mathcal{E}}
\newcommand{\Rcal}{\mathcal{R}}
% changed here
\newcommand{\TTop}{\mathcal{T}}
% end changed here
\newcommand{\WinningSet}[1]{\Ecal(C, n, \Rcal, #1)}
\newcommand{\WinningSetR}[2]{\Ecal(C, n, #1, #2)}
\newcommand{\wsonea}{\Ecal(C, n, R_1, a)}
\newcommand{\Wsp}{\Ecal^{sp}_n}
\newcommand{\todo}[1]{\textcolor{blue}{[TODO: #1]}}

\newcommand{\bicatb}{\mathcal{B}}

\newcommand{\catc}{\mathcal{C}}
\newcommand{\catcop}{\catc^{op}}
\newcommand{\catd}{\mathcal{D}}
\newcommand{\cate}{\mathcal{E}}
\newcommand{\eqdef}{\triangleq}
\newcommand{\id}{\mathbf{id}}
\newcommand{\idbicat}{\mathbf{Id}}
\newcommand{\Id}{\mathbf{Id}}
\newcommand{\righttop}[1]{\xrightarrow{#1}}
\newcommand{\lefttop}[1]{\xleftarrow{#1}}
\newcommand{\Ob}{\mathbf{Ob} \ }
\newcommand{\Set}{\mathbf{Set}}
\newcommand{\SpanC}{\mathbf{Span(\catc)}}
\newcommand{\SpanSet}{\mathbf{Span(Set)}}
\newcommand{\Cat}{\mathbf{Cat}}
\newcommand{\SetMC}{\Set^\times}
\newcommand{\spancat}[5]{#1 \lefttop{#2} #3 \righttop{#4} #5}
\newcommand{\cospancat}[5]{#1 \righttop{#2} #3\lefttop{#4} #5}
\newcommand{\pb}[5]{#1 \times_{#2#3#4} #5}
\newcommand{\pbs}[3]{#1 \times_{#2} #3}
\newcommand{\pbup}[3]{\langle #1, #2 \rangle_{#3}}
\newcommand{\pbm}[3]{#1\times_{#2}#3}
% \newcommand{\Tree}{\mathcal{T}}
\newcommand{\npbup}[2]{\langle #1 \rangle_{#2}}
\newcommand{\npbuptmp}[1]{\llbracket #1 \rrbracket}
\newcommand{\Scal}{\mathcal{S}}
\newcommand{\CtxZero}{\mathbf{Ctx_0}}
\newcommand{\F}{\mathbb{F}}
\newcommand{\Fop}{\mathbb{F}^{op}}
\newcommand{\CtxOne}{\mathbf{Ctx_1}}
\newcommand{\LT}{\mathbb{L}}
\newcommand{\catv}{\mathcal{V}}
\newcommand{\struct}{\mathbf{struct}}
\newcommand{\arity}[2]{( #1 \ | \ #2)}
\newcommand{\FinProd}{\mathsf{FinProd}}
\newcommand{\D}{\mathsf{D}}
% \newcommand{\yo}{\textmc{よ}}
\newcommand{\yo}{\textbf{yo}}
\newcommand{\Bij}{\mathbf{Bij}}
\newcommand{\Inj}{\mathbf{Inj}}
\newcommand{\FinSet}{\mathbf{FinSet}}
\newcommand{\BijSystem}{\mathsf{Bij}}
\newcommand{\InjSystem}{\mathsf{Inj}}
\newcommand{\CartSystem}{\mathsf{Cart}}
\newcommand{\dayconv}{\otimes_{\text{Day}}}
\newcommand{\Law}{\mathbf{Law}}
\newcommand{\LawEnriched}{\mathbf{Law}^{\text{enr}}}

\newcommand{\terminalmorph}{\langle\rangle}
\newcommand{\prodmorph}[2]{\langle#1, #2\rangle}

% \newcommand{\ul}[1]{\underline{#1}}
\newcommand{\yozero}{\yo(0)}
\newcommand{\yoone}{\yo(1)}
\newcommand{\yon}{\yo(n)}
\newcommand{\yom}{\yo({m})}
\newcommand{\yop}{\yo(p)}
\newcommand{\op}{\mathbf{op}}
\newcommand{\Ops}{\mathcal{O}}
\newcommand{\typing}[3]{#1 \ | \ #2 \vdash #3}
\newcommand{\typingfilled}[1]{\Gamma \ | \ \Delta \vdash #1}
% \newcommand\eye{\ensurestackMath{\stackinset{c}{}{c}{-.33pt}%
%   {\bullet}{\bigcirc}}}
% \newcommand\smallereye{\scalerel*{\eye}{x}}
% \newcommand{\eyemid}{ \ \eye \ }
\newcommand{\eye}{\odot}
\newcommand{\eyemid}{\odot}
\newcommand{\sem}[1]{\llbracket#1\rrbracket}
\newcommand{\metasub}[2]{\{#1/#2\}}
\newcommand{\inl}{\mathbf{inl}}
\newcommand{\inr}{\mathbf{inr}}

% I/O Quantum things
\newcommand{\quantum}{\mathsf{Quantum}}
\newcommand{\io}{\mathsf{I/O}}
\newcommand{\iotq}{\io \otimes \quantum}

% Useful Categories
\newcommand{\CPU}{\mathbf{CPU}}
\newcommand{\CPTP}{\mathbf{CPTP}}
\newcommand{\CP}{\mathbf{CP}}

% Operations
\newcommand{\new}{\mathsf{new}}
\newcommand{\apply}{\mathsf{apply}}
\newcommand{\applyU}{\mathsf{apply}_U}
% \newcommand{\measure}{\mathsf{measure}}
\newcommand{\opmeasure}{\mathsf{measure}}
\newcommand{\discard}{\mathsf{discard}}

% I/O operations
\newcommand{\opin}{\mathsf{in}}
\newcommand{\opout}{\mathsf{out}}
\newcommand{\opcin}{\mathsf{cin}}
\newcommand{\opcout}{\mathsf{cout}}
\newcommand{\opskip}{\mathsf{skip}}

% Vectors
\newcommand{\veca}{\Vec{a}}
\newcommand{\vecb}{\Vec{b}}
\newcommand{\vecc}{\Vec{c}}
\newcommand{\vecd}{\Vec{d}}
\newcommand{\vecx}{\Vec{x}}
\newcommand{\vecy}{\Vec{y}}
\newcommand{\vecz}{\Vec{z}}
\newcommand{\vecp}{\Vec{p}}
\newcommand{\vecq}{\Vec{q}}


\newcommand{\statein}{\mathbf{in}}
\newcommand{\traceset}{\{\opin, \opout\}^*}
\newcommand{\Stream}{\mathsf{Stream}}

\newcommand{\Mat}{\mathcal{M}}
\newcommand{\qubit}{\Mat_2(\complex)}
\newcommand{\MnC}{\Mat_n(\complex)}
\newcommand{\MkC}{\Mat_k(\complex)}
\newcommand{\MtnC}{\Mat_{2^n}(\complex)}

\newcommand{\qbit}{\mathbf{qbit}}
\newcommand{\qbitn}[1]{\qbit^{\otimes #1}}
\newcommand{\bit}{\mathbf{bit}}

\newcommand{\rhotilde}{\tilde{\rho}}
\newcommand{\psitilde}{\tilde{\psi}}
\newcommand{\Config}{\mathbf{Config}}
\newcommand{\Dist}{\mathcal{D}}
\newcommand{\config}[3]{\left\langle #1, #2, #3 \right\rangle}
\newcommand{\trace}{\mathbf{tr}}
\newcommand{\unifiedconfig}[3]{\left[ #1, #2, #3 \right]}

\newcommand{\probred}[2]{\downarrow_{#2}^{#1}}
\newcommand{\probredp}[1]{\downarrow_{#1}^{t}}
\newcommand{\probredfilled}{\downarrow_{p}^{t}}
\newcommand{\probequiv}{\simeq_{\text{pr}}}
\newcommand{\psequiv}{\simeq_{pr}^{\opskip}}


\newcommand{\qured}[3]{\Downarrow_{#3}^{#1, #2}}
\newcommand{\quredp}[1]{\Downarrow_{#1}^{\Gamma, t}}
\newcommand{\quredfilled}{\Downarrow_{p}^{\Gamma, t}}
\newcommand{\quequiv}{\simeq_{qu}}

\newcommand{\alphaequiv}{\simeq_\alpha}

\newcommand{\phibar}{\bar{\phi}}
\newcommand{\tostar}{\to^*}

\newcommand{\LHS}{\mathsf{LHS}}
\newcommand{\RHS}{\mathsf{RHS}}

\newcommand{\Path}[1]{\mathsf{Path}\left(#1\right)}

% Chapter 2 algebraic theories
\newcommand{\Mon}{\mathsf{Mon}}

\newcommand{\Tree}{\mathsf{Tree}}

\newcommand{\return}[1]{\mathbf{return} \ #1}
\newcommand{\Kl}[1]{\mathsf{Kl}_{#1}}

% chapter 5
\newcommand{\CPinf}{\CP^\infty}
\newcommand{\Mod}{\mathsf{Mod}}
\newcommand{\Aux}{\mathsf{Aux}}
\newcommand{\qbitin}{\qbit^{\opin}}
\newcommand{\qbitout}{\qbit^{\opout}}
\newcommand{\CPM}{\mathbf{CPM}}
\newcommand{\CPMbar}{\overline{\CPM}}
\newcommand{\geqlowner}{\sqsupseteq}
\newcommand{\leqlowner}{\sqsubseteq}
\newcommand{\lub}{\bigsqcup}
\newcommand{\inj}{\mathsf{inj}}
\newcommand{\projsf}{\mathsf{proj}}
\newcommand{\Run}{\mathsf{Run}}
\newcommand{\Connect}{\mathsf{Connect}}

\newcommand{\ICal}{\mathcal{I}}


\begin{document}

\title{Semantics and Equational Axiomatisation for Quantum Communication: Presentation Proposal}

\author{Theo Wang}
\authornote{This presentation is based on my MSc thesis \cite{wang_algebraic_2024} with Sam Staton (Oxford). }
\email{theo.wang@spc.ox.ac.uk}
\affiliation{%
  \institution{University of Oxford}
  \city{Oxford}
  \country{United Kingdom}}

\maketitle

\section{Overview of the Presentation}

\paragraph{What.} We systematically study the semantics of quantum programs extended with classically controlled quantum communication:  sending and receiving qubits, possibly depending on a classical measurement outcome. For example, the following is naturally expressed in imperative syntax, but not in the standard quantum circuit model.
\begin{lstlisting}
  p: qbit = in(); q: qbit = new(|0>); 
  if (measure(p) == |1>) { out(q); }
\end{lstlisting}
The end goal would be to develop an axiomatisation of program equivalence between communicating quantum programs. 

\paragraph{Why.} Quantum communication is an important programming primitive in distributive quantum computing~\cite{DQC_survey_Caleffi_2024}. Textbook quantum protocols such as teleportation \cite{Teleportation_PhysRevLett.70.1895} and BB84 \cite{BB84_BENNETT20147} already implicitly use communication primitives. Prior work already treats quantum communication within quantum process calculi~\cite{qCCS_ying2010algebraquantumprocesses,CQP_gay2004communicatingquantumprocesses,QPAlg_lalire2004processalgebraicapproachconcurrent, lqCCS_10.1145/3632885}, where communication is entangled with the additional complexity of concurrency.

\paragraph{How.} We study the individual operations of quantum communication \emph{in isolation}, as algebraic effects, following the tradition of Plotkin and Power~\cite{plotkin_power_2002}. We do so along three axes, following the triangle summarised by \cref{fig:triangle}. We propose a new algebraic theory for quantum communication building on \cite{statonpopl15}, and study an operational and a denotational model of the theory. Finally, we show that denotational and operational equalities coincide.  

\begin{figure}[ht]
  \centering
  \begin{tikzpicture}[scale=0.7, every node/.style={align=center, font=\footnotesize}]
    \node (A) at (0, 1.2) {{\bfseries Algebraic theory} $\iotq$};
    \node (B) at (-3.4, -1.2) {{\bfseries Operational semantics} \\ quantum streams};
    \node (C) at (3.4, -1.2) {{\bfseries Denotational semantics} \\ free I/O monad on $\CPinf$};
    \draw[->] (A) -- node[midway, above left, font=\scriptsize] {\cref{thm:soundopsem}} (B);
    \draw[->] (A) -- node[midway, above right, font=\scriptsize] {\cref{thm:sounddensem}} (C);
    \draw[<->] (B) -- node[midway, below, font=\scriptsize] {\cref{thm:adeqfullabs}} (C);
  \end{tikzpicture}
  \caption{Three semantics of quantum communication and how they relate.}
  \label{fig:triangle}
\end{figure}

\section{Background: Axiomatising Quantum Computing}
We already know how to build the same triangle for qubit quantum computing without communication. Our starting point is \cite{statonpopl15}'s axiomatisation for qubit quantum computing, $\quantum$, based on \textit{parameterised} algebraic theories (PAT). PATs are a generic framework for axiomatising the behaviour of computations dependending on certain resources in terms of operations and equations. Here, resources will correspond to qubits.

\paragraph{Operations} $\op: \arity{p}{m_1...m_k}$ represent computations which take $p$ qubits, classically choose one of the $k$ branches, say, $i$, and return $m_i$ qubits. $\quantum$ has three operations: state preparation $\new: (0 \mid 1)$, unitary evolution $\applyU: (n \mid n)$ for each $2^n \times 2^n$ unitary $U$, and classical control $\opmeasure: (1 \mid 0, 0)$.

\paragraph{Terms} generated from the operations, follow the same intuition: a term $\typing{x_1: m_1...x_k: m_k}{q_1...q_p}{c}$ takes $p$ input qubits $q_1...q_p$ and calls one of the $k$ continuations $x_1...x_k$ with the corresponding number of qubits. We can thus write for example the program 
  \[
  \typing{x: 1, y: 1}{q_1, q_2}{\opmeasure(q_1, x(q_2), y(q_2))}
  \]
  which measures a $q_1$ and continues with $x$ or $y$ with $q_2$ based on the measurement outcome. 
\paragraph{Equations} are given in the form $\typing{\Gamma}{\Delta}{c=c'}$. $\quantum$ for example specifies effect of measuring a $\ket|0>$ qubit using the following equation:
\[
\typing{x: 0, y: 0}{\cdot}{\new(q. \opmeasure(q, x, y)) = x}
\]

\paragraph{Models: Denotational and Operational} \cite{statonpopl15} proceeds to define a denotational model for the $\quantum$ theory, where each term $\typing{x_1:m_1...x_k:m_k}{q_1...q_p}{c}$ is interpreted as a complete positive unital (CPU) map $$\sem{c}: \bigoplus_{i=1...k}\Mat_i(\complex) \to \Mat_p(\complex).$$ Remarkably, this model is \emph{fully complete}: every CPU map of this type is the interpretation of some term $c$, and two CPU maps are equal iff their term representation are derivably equal in the equational theory. 

We can in fact give an equivalent operational model for $\quantum$. Briefly, configurations would look like $\langle \rho, c \rangle$ where $\rho$ is the program state, $c$ the running program, and where transitions $\to_p$ occur probabilistically. We can then formulate in terms of $\to_p$ an operational notion of program equivalence $\quequiv$ equivalent with denotational equality.

\section{Semantics for Quantum Communication}
We proceed to extend our language with communication primitives, and strive to recover the same triangle. 

\subsection{An Algebraic Theory for Quantum Communication}
Extending $\quantum$ with two operations: $\opin: (0 \mid 1)$ and $\opout: (1 \mid 0)$ lets us represent classically controlled communication. We can express the example program as $\typing{x: 0}{\cdot}{\opin(p. \new(q. \opmeasure(p, x, \opout(q, x))))}$. We further add equations to make $\opin$ and $\opout$ commute with every $\quantum$ operation. What we now obtain is $\iotq$, where $\otimes$ denotes the commutative combination of two theories. This choice is motivated by quantum theory. For instance,  
$$\opmeasure(a,\; \opout(b, x),\; \opout(b, y)) = \opout(b,\; \opmeasure(a, x, y))$$
is required to derive the equivalence between: preparing a Bell pair, sending one end and discarding the other end; and tossing a fair coin and sending the outcome as a qubit. 

\subsection{An Operational Model}
The notion of communication that our theory soundly axiomatises can be demonstrated through an operational model. We work in, $\CP$ the category of finite dimensional C* algebras and complete positive maps, constructed as in \cite{selinger_dagger_compact_closed_cat_cp_maps}. 

The operational semantics is based on a \emph{quantum stream}, specifying the environment our program communicates with. The idea is then to have a global quantum state $\rho: \complex \to S \otimes \qbit \otimes H$, where $S$ is the program state, $H$ is the hidden stream state, and $\qbit$ is the next input qubit. Then, a stream is a state machine reacting to the program's communications: formally, a stream is specified by a hidden state type $H_\sigma$ and transition maps $T_{\opin}(\sigma)$, $T_{\opout}(\sigma)$ for every communication trace $\sigma \in \traceset$. 

We can then build an operational semantics with configuration $\config{\rho}{\sigma}{c}$, where $\rho$ is the global state, $\sigma \in \{\opin, \opout\}^*$, $c$ is the program. Given a stream specification $t$ we can give a small step transition relation $\to^t_p$ defined inductively on the program. For example for the output operation, the rule is as in \cref{fig:opsemout}: the output is performed, and the stream state is updated according to $t$. It is then not difficult to extend $\quequiv$ to communicating quantum programs, to obtain a sound model of $\iotq$: 
\begin{theorem}
  \label{thm:soundopsem}
  $\typing{\Gamma}{\Delta}{c=c'} \implies \typing{\Gamma}{\Delta}{c \quequiv c'}$.
\end{theorem}
\begin{figure}
  \begin{equation}
  \config{\tikzfig{meas_out_init}}{\sigma}{\opout(q, c)} \to_1 \config{\tikzfig{out_post}}{\sigma + \opout}{c} \tag{\text{out}}
\end{equation}
\caption{Operational Rule for $\opout$}
\label{fig:opsemout}
\end{figure}

\subsection{A Denotational Model}
\paragraph{Strategy}

\paragraph{The Quantum I/O Monad}
\paragraph{Denotational Semantics}

\begin{theorem}
  \label{thm:adeqfullabs}
  $\typing{\Gamma}{\Delta}{c \quequiv c'} \iff \sem{c} = \sem{c'}$.
\end{theorem}
\begin{corollary}
  \label{thm:sounddensem}
  $\typing{\Gamma}{\Delta}{c=c'} \implies \sem{c} = \sem{c'}$.
\end{corollary}

\section{Limitations and Future Work}

We have thus obtained a sound axiomatisation of two notions of quantum communication, one operationally presented, one denotationally presented, which end up being equivalent. However, unlike the quantum computing case, we do not have completeness; it is unclear whether two programs $c$ and $c'$ being denotationally equal implies that they are derivably equal in $\iotq$. Achieving this is our immediate next step. 

\section{Background: parameterised algebraic theories and QUANTUM}

A parameterised algebraic theory (PAT)~\cite{statonpopl15} is a signature of operations together with equations over their terms, whose operations carry \emph{linear parameter contexts}: a typing judgement $\typing{\Gamma}{\Delta}{c}$ has $\Delta$ a list of parameter names (here, qubits) and $\Gamma$ a list of continuation variables with their valences. Operations $\op:\arity{p}{m_1\dots m_k}$ consume $p$ parameters and bind $m_i$ fresh parameters in each of $k$ continuations. Read as a type discipline, this is a linear type system for resource-sensitive effects.

Two combinators are standard: the \emph{sum} (the free combination, no equations imposed between the two signatures) and the \emph{tensor} (which adds the equation that every operation of one theory commutes with every operation of the other). Their constructions for the PATs of~\cite{statonpopl15} are worked out in~\cite{wang_algebraic_2024}.

Staton's $\quantum$~\cite{statonpopl15} is a PAT for qubit quantum computation with operations $\new:\arity{0}{1}$ (state preparation), $\applyU:\arity{n}{n}$ (unitary evolution) and $\opmeasure:\arity{1}{0,0}$ (computational-basis measurement with classical control), together with quantum-theoretic equations such as the principle of deferred measurement. It is sound and complete with respect to its standard operator-algebra interpretation in $\CP$, the finite-biproduct completion of the category $\CPM$ of matrix algebras and CP maps. The analogous picture to Figure~\ref{fig:triangle} is settled in this I/O-free case. We extend it to communication.

\section{The algebraic theory for I/O and QUANTUM}

Extending $\quantum$ with $\opin:\arity{0}{1}$ (receive a qubit) and $\opout:\arity{1}{0}$ (send a qubit) lets us write classically controlled communication directly. The theory $\iotq$ is the \emph{tensor} (not sum) of $\quantum$ with a standard $\io$ theory~\cite{sums_tensors_HYLAND200670}. There are two design points. First, $\opin$ and $\opout$ do not commute with each other --- the order of I/O events is part of the observable behaviour of a program. Second, $\opin$ and $\opout$ do commute with all of $\quantum$'s operations; this is motivated by quantum theory, since for instance
$$\opmeasure(a,\; \opout(b, x),\; \opout(b, y)) = \opout(b,\; \opmeasure(a, x, y))$$
is needed to make a familiar quantum equivalence --- preparing a Bell pair, sending one qubit and discarding the other equals tossing a fair coin and preparing the outcome --- provable in the theory~(\cite{wang_algebraic_2024}, eq.~4.1.4). Tensors of theories can collapse (the Eckmann--Hilton argument for unital binary operations is the textbook example); whether $\iotq$ collapses is ruled out \emph{a posteriori} by the two models in \S\S4--5.

\section{Operational semantics: quantum streams}

A \emph{quantum stream} is a $\traceset$-indexed family $q:\traceset \to \Ob\CP$ together with CPTP transitions $T(\sigma,\opin)$ and $T(\sigma,\opout)$ that update a hidden environment state at each I/O event. Streams are coalgebraic in flavour; we comment on a coalgebraic recasting in \S7. A configuration $\config{\rho}{\sigma}{c}$ pairs a normalised global state $\rho$ (over program qubits, communication qubit, and hidden state), an I/O trace $\sigma$, and a program $c$. The small-step probabilistic transition $\to^t_p$ is defined inductively in the program structure: measurement reduces probabilistically; other operations reduce deterministically; a continuation call terminates.

The induced \emph{quantum equivalence} $c_1 \quequiv c_2$ holds when, for every stream and every initial state, the two programs yield the same probability-weighted distribution over (final state, continuation called, final I/O trace) triples. Terms of $\iotq$ modulo $\quequiv$ form a model of $\iotq$~(\cite{wang_algebraic_2024}, Thm.~4.3.25): the equations of the PAT are validated by the stream dynamics.

\section{Denotational semantics: a free I/O monad}

The operational semantics lives in $\CP$, which has only finite biproducts; the denotational sum over all I/O traces below needs infinitely many. Following~\cite{PaganiSV13}, $\CPinf$ is the infinite-biproduct, dcpo-enriched completion of $\CPM$, inheriting compact closure from $\CPM$. On $\CPinf$ we define
$$
  T(X) \;\eqdef\; \mu Y.\; \bigl(\qbitin \otimes Y\bigr) \oplus \bigl(\qbitout \otimes Y\bigr) \oplus X \;\;\cong\;\; \bigoplus_{\sigma \in \traceset} \Aux(\sigma) \otimes X,
$$
where $\Aux(\epsilon)=\complex$ and, for $s\in\{\opin,\opout\}$, $\Aux(s\!\cdot\!\sigma)=\qbit^{s}\otimes\Aux(\sigma)$; the superscripts $\opin/\opout$ are bookkeeping, both $\qbitin$ and $\qbitout$ are just $\qbit$. For every I/O trace $\sigma$, the corresponding summand tensors the qubits exchanged along $\sigma$ with a value of type $X$.

The input case deserves a comment. An input operation morally has type $\qbit\!\multimap\!Y$. Compact closure of $\CPinf$ gives $\qbit\!\multimap\!Y \cong \qbit^{*}\otimes Y$, and the further fact that $\qbit$ is self-dual in $\CPM$ gives $\qbit^{*}\cong\qbit$, so the input summand can be written as $\qbitin \otimes Y$. The two readings are reconciled by a single string-diagrammatic move: the unnormalised maximally entangled state acts as a cup, bending the input wire.

Interpreting parameter contexts as tensor powers of $\qbit$ and continuation contexts as direct sums, each term $\typing{\Gamma}{\Delta}{c}$ denotes a Kleisli morphism
$$\sem{c} \;:\; \sem{\Delta} \longrightarrow T\sem{\Gamma} \;\cong\; \bigoplus_{\sigma \in \traceset,\; (x:n) \in \Gamma} \Aux(\sigma) \otimes \qbitn{n},$$
a family of CP maps indexed by (final I/O trace, continuation called) --- the same indexing already used by $\quequiv$. The object $\complex$ is a model of $\iotq$ in $\Kl{T}^{\mathrm{op}}$~(\cite{wang_algebraic_2024}, Cor.~5.4.2).

\section{Adequacy and full abstraction}

The two concrete models coincide on program equivalence.
\begin{theorem}[\cite{wang_algebraic_2024}, Thm.~5.3.4 (adequacy) and Thm.~5.3.6 (full abstraction)]\label{thm:adequacyfullabs}
  For all $\typing{\Gamma}{\Delta}{c_1, c_2}$:\quad
  $\sem{c_1}=\sem{c_2}$\quad iff\quad $c_1 \quequiv c_2$.
\end{theorem}
This is the kind of abstract--concrete correspondence the workshop addresses: both sides are independent concrete techniques, and Theorem~\ref{thm:adequacyfullabs} shows they agree, with the abstract theory $\iotq$ as the common interface that each refines. The proof goes via a \emph{sum-of-paths} characterisation matching the operational sum over reduction sequences to the categorical sum over biproduct components, term by term. The matching of Theorem~\ref{thm:adequacyfullabs} does not, however, settle \emph{completeness} of $\iotq$: whether the equational theory derives every identification validated by the two concrete models is open~(\cite{wang_algebraic_2024}, Conjecture~5.4.4).

\section{Discussion and open directions}

The pattern \emph{algebraic theory + two concrete models + adequacy/FA} echoes themes in the algebraic-effects programme of Plotkin and Power~\cite{plotkin_power_2002} and is instantiated, for $\quantum$ alone, in~\cite{statonpopl15}. The quantum I/O case is technically non-trivial because compact closure, infinite biproducts, and the possibility of a tensor-of-theories collapse each constrain the construction.

We would welcome discussion at the workshop on three open directions. \emph{(a)} \textbf{Completeness} of $\iotq$ for $\quequiv$: what would a suitable axiomatic logic look like for a matching like Theorem~\ref{thm:adequacyfullabs}? \emph{(b)} \textbf{Coalgebraic reading.} The stream of \S4 is coalgebraic, and Ceragioli et al.'s coalgebraic model of quantum bisimulation~\cite{ceragiolicoalgebraic} appears to provide a coalgebraic framing of related phenomena; we are also curious how our denotational semantics relates to the $\mathsf{Caus}$ construction of Kissinger--Uijlen~\cite{kissinger2019categorical}. \emph{(c)} \textbf{Algorithmic verification.} The denotational semantics is probabilistic and indexed by I/O traces; is there a decidable fragment, or a probabilistic / quantum Hennessy--Milner logic against which equivalence in $\iotq$ could be model-checked?

\bibliographystyle{ACM-Reference-Format}
\bibliography{refs}

\end{document}
